Supervised training of a finite-element surrogate requires a full solve for every training
example, the very cost the surrogate should remove. For linear elastostatics we train on
the assembled discrete potential energy instead. Minimising it is exactly supervised
regression in the stiffness norm, which measures the strain energy of the error, so the
surrogate learns from a weighted mean-square stress error without solving any example. In pre-registered comparisons on three-dimensional tetrahedral
meshes, a transformer trained on the energy alone was more accurate, on every metric,
than the same network trained on the solutions of the same instances. A supervised graph
network had smaller displacement errors on most instances but larger stress-field errors
on almost all: displacement error ranked the two in the opposite order to stress error.
Every supervised network there returned some predictions worse than the zero field in the
energy norm, failures the energy detects without reference solutions. In two dimensions
the label-free transformer had a larger displacement error than its supervised
counterpart but smaller stress errors; trained on the stiffness-norm error instead, with
the same labels, the counterpart halved its von Mises error, as pre-registered, and its
errors lay in softer stiffness eigenmodes, as measured afterwards: with the labels fixed,
the norm of the loss decided the stress accuracy. No network transferred to a
substantially finer three-dimensional mesh without retraining; a post hoc, closed-form,
label-free rescaling removed most of the label-free network's error there. The discrete
energy thus serves as training objective, evaluation metric and failure detector.
